\documentclass[10pt,a4paper]{amsart}
\usepackage[margin=19mm]{geometry}
\usepackage{amsmath,amssymb,mathtools}
\usepackage[scr=rsfs,cal=euler]{mathalfa}
\usepackage{xfrac}
\usepackage[unicode,hidelinks,bookmarks=false]{hyperref}
\newcommand{\ie}{i.e.\ }
\newcommand{\cf}{cf.\ }
\newcommand{\resp}{respectively\ }
\newcommand{\Subsection}{\subsection}
\providecommand{\nobreakdash}{\nobreak\hbox{-}}
\setlength{\emergencystretch}{2em}
\multlinegap0pt
\usepackage{amssymb}
\usepackage{xcolor}
\newtheorem{theorem}{Theorem}
\newtheorem{lemma}{Lemma}
\newcommand{\cA}{{\mathcal{A}}}
\title{Explicit Burgess inequalities for cubefree moduli}
\author{Elchin Hasanalizade, Hua Lin, Andradis Luna Mart\'inez, Greg Martin, Enrique Trevi{\~n}o}
\date{Source: arXiv:2511.17778v2 (CC BY 4.0)}
\begin{document}
\maketitle

\noindent\emph{Source and licence.} Explicit Burgess inequalities for cubefree moduli. Version arXiv:2511.17778v2 (CC BY 4.0), distributed by arXiv under the Creative Commons Attribution 4.0 International licence, http://creativecommons.org/licenses/by/4.0/, as recorded in the arXiv licence metadata for this exact version and shown on the abstract page https://arxiv.org/abs/2511.17778v2. Full licence notice: \texttt{licenses/arxiv-2511.17778-CC-BY-4.0.txt}.\par\smallskip

\noindent\textit{Editorial scope.} Counted unit: the article's theorem 1 (source0.tex lines 106--112). The article's complete original proof of it is delivered with it (source0.tex lines 875--925, one \texttt{proof} environment of 51 lines, which the article places away from the statement). No mathematics is written, repaired or re-derived: the statement and the proof are the article's own lines, in the article's own notation, and no retained line is substituted or re-flowed.  Retained uncounted as the source-local context the proof needs: lines 114--131; lines 135--138; lines 304--307; lines 494--500; lines 536--541; lines 565--575; lines 624--637; lines 660--662; lines 664--666; lines 671--684; lines 757--759; lines 798--806; lines 828--830; lines 834--842.  Nothing was omitted from any retained span, so the package carries no omission record.  No retained line contains a citation, so the wrapper carries no bibliography and no bibliography entry.  No retained line contains a figure environment, an image-inclusion command or drawing code, so no image is delivered and the package declares no asset dependency.  Editorial formatting only, in addition: an AMS \texttt{amsart} preamble that restates the article's own declarations and macros used by the retained lines, namely the environments \texttt{theorem}, \texttt{lemma} on separate counters, together with the standard packages those lines need.  The retained environments print in this document as Theorem 1, Lemma 1 in order of appearance, whereas the article prints the same items with the numbering of its own full document.  The complete native source of the article as distributed in the arXiv e-print for this exact version is bundled unchanged as \texttt{sources/189610/source0.tex}.
\begin{table}[h!]
\begin{center}
  \begin{tabular}{ | c | c | c|c|}
    \hline
    $r$ & $C(r)$ & $D(r)$ & $C^*(r)$ \\ \hline
    2 & 14.275 & 7.844 &  12.364\\ \hline
    3 & 5.135 & 4.390 & 4.849\\ \hline
    4 & 3.556 & 3.289 & 3.448\\ \hline
    5 & 2.878 & 2.740 & 2.823\\ \hline
    6 & 2.496 & 2.410 & 2.463\\ \hline
    7 & 2.248 & 2.189 & 2.227\\ \hline
    8 & 2.075 & 2.030 & 2.060\\ \hline
    9 & 1.945 & 1.911 & 1.934\\ \hline
    $\ge$10 & 1.846 & 1.818 & 1.837\\ \hline
  \end{tabular}
  \caption{Constants in the inequalities for values of $r$ in Theorems \ref{main thm 1} and \ref{main theorem 2}.}\label{Table 1}
\end{center}
\end{table}

\begin{theorem}\label{main theorem 2}
    Let $\chi$ be a primitive Dirichlet character modulo $q$. Let $C^{*}(r)$ be defined as in Table \ref{Table 1}. Then, for $q \ge \max\{10^{1143}, 2^{4r-2}\}$, if $r=2$ or $q$ is cubefree, we have
    $$|S_{\chi}(M,N)| \le C^{*}(r)N^{1-\frac{1}{r}} q^{\frac{r+1}{4r^2}}(\log{q})^{\frac{1}{2r}}\left((4r)^{\omega(q)}m_r(q)\right)^{\frac{1}{2r}}\left(\frac{q}{\phi(q)}\right)^{\frac{1}{r}}.$$
\end{theorem}

\begin{equation} \label{eq desired bound}
| S_\chi(M,N) | \le E_{q,r}(N), \quad\text{where }
E_{q,r}(N) = C N^{1-1/r} q^{\frac{r+1}{4r^2}} (\log q)^{\frac{1}{2r}} T(q),
\end{equation}

\begin{lemma} \label{B will be at least 10 if it kills us lemma}
For $r\ge 2$, let $a(r) = 2\left(3.0758r+1.38402\log(4r)-1.5379\right)\log{2}$. If $q \ge e^{e^{a(r)}}$, then $$B_1(r,q)\ge 2^{2-2/r}r^2\left(\frac{1-\frac{1}{r}}{r!}\right)^{\frac{1}{r}},$$
where 
\begin{equation}\label{def B_1}
B_1(r,q) = r^2q^{\frac{1}{2r}}\biggl( \frac{r-1}{r!2r(4r)^{\omega(q)}m_r(q)} \biggr)^{1/r}.
\end{equation}
\end{lemma}

\begin{lemma} \label{B will be at least 10 if it kills us lemma 2}
Let $r\ge 2$. If $q \ge 2^{4r-2}$, then $$B_2(r,q)\ge 2^{2-2/r}r^2\left(\frac{1-\frac{1}{r}}{r!}\right)^{\frac{1}{r}},$$
where 
\begin{equation}\label{def B_2}B_2(r,q) = r^2q^{\frac{1}{2r}}\biggl( \frac{r-1}{r!2r} \biggr)^{1/r}.
\end{equation}
\end{lemma}

\begin{lemma} \label{A lower bound lemma}
Let $N \ge q^{\frac{1}{4}+\frac{1}{4r}}$ and $r\ge 2$ be positive integers. Suppose 
$$A_i(r,q) = \frac{\kappa N}{B_i(r,q)},$$
for $i=1,2$ where $B_1(r,q)$ is defined as in \eqref{def B_1}, $B_2(r,q)$ as in \eqref{def B_2}, and let
\begin{equation} \label{choice of kappa}
\kappa = \biggl( \frac{(B_i(r,q)-1) B_i(r,q)^{1-1/r} }{2 r (B_i(r,q)+1)^{2-1/r}}
\biggr)^{r/(r-1)}.
\end{equation}
Let $t_1(r) = e^{e^{a(r)}}$, where $a(r)$ is defined as in Lemma \ref{B will be at least 10 if it kills us lemma} and $t_2(r) = 2^{4r-2}$. Then, for $q\ge \max\bigl\{ 10^{1143}, (10r)^{18}, t_i(r) \bigr\}$, we have 
$$A_i(r,q) \ge \max\biggl\{31, \frac{2^{\omega(q)}q}{\phi(q)}\biggl\}.$$
\end{lemma}

\begin{lemma}\label{lem: alpha bound}
Let $r,q\ge 2$ and $N \ge q^{\frac{1}{4}+\frac{1}{4r}} $ be integers and let $i \in \{1,2\}$. Let $A = A_i(r,q)$, $B = B_i(r,q)$, $\kappa$ be defined as in \eqref{choice of kappa}, and $t_i(r)$ be defined as in Lemma \ref{A lower bound lemma}.
Let
\begin{equation*}
    \alpha
    =
    \left[ \frac{\left(\kappa\left(\frac{\phi(q)}{q}\right)+\frac{B2^{\omega(q)-1}}{N} \right)^2}{B}
    +
    2\kappa\left(\frac{6}{\pi^2}\log{(A)}+\delta+\frac{\log{(A)}+1}{3A}\right)\right],
\end{equation*}
where $\delta = \frac{6\gamma}{\pi^2} - \frac{36\zeta'(2)}{\pi^4} \approx 0.6974 $.
Then, for $q\ge \max\{10^{1143}, (10r)^{18}, t_i(q)\}$,
$$\alpha \le \frac{32}{31}\kappa \log{q}.$$
\end{lemma}

\begin{align}\label{eq T1}
    T_1(q) = \bigl((4r)^{\omega(q)} m_r(q)\bigr)^{\left({\frac{1}{2r}-\frac{1}{2r^2}}\right)} \biggl( \frac q{\phi(q)} \biggr)^{1/r},
\end{align}

\begin{align}\label{eq T2}
    T_2(q) = \bigl((4r)^{\omega(q)} m_r(q)\bigr)^{\left({\frac{1}{2r}}\right)} \biggl( \frac q{\phi(q)} \biggr)^{1/r}.
\end{align}

\begin{lemma}\label{lemma to set up the main theorems}
For $i\in\{1, 2\}$ and $T_i$ be defined as in \eqref{eq T1} and \eqref{eq T2}. Let $s, \alpha, u, $ and $w$ be defined as in \eqref{eq s}, \eqref{eq alpha}, \eqref{eq u}, and \eqref{w}, respectively. Then
    \begin{equation*}
\begin{split}
    \left| S_\chi(M, N) \right|
    \le
    \left(\frac{B}{B-1}\right)N^{1-1/r}q^{(r+1)/{4r^2}}(\log q)^{\frac{1}{2r}}
    \left(
    u(\alpha w)^{\frac{1}{2r}}
    +sT_i(q)
    \right).
\end{split} 
\end{equation*}
\end{lemma}

\begin{equation}\label{eq s}
   s=\frac{2C\kappa^{1-1/r}}{2-1/r}\left(\frac{B+1}{B}\right)^{2-1/r},
\end{equation}

\begin{equation}\label{eq alpha}
    \alpha
    =
    PB/N^2
    =
    \left[ \frac{\left(\kappa\phi^*+\frac{B2^{\omega(q)-1}}{N} \right)^2}{B}
    +
    2\kappa\left(\frac{6}{\pi^2}\log{(A)}+\delta+\frac{\log{(A)}+1}{3A}\right)\right].
\end{equation}

\begin{equation}\label{eq u}
u=\frac{\left(N/\#\cA\right)^{1/{r}}}{q^{\frac{1}{2r^2}}},
\end{equation}

\begin{equation}
\label{w}
    w
    = \frac{f_i}{r^2 \beta^{r+1} r! \log{q}}.
   % \frac{2(4r)^{\omega(q)}m_r(q)}{\beta r(r-1)\log q} &=
    %\frac{2(4r)^{\omega(q)}m_r(q)}{r(r-1)\log q} \left(\frac{r-1}{r!2r\left(4r\right)^{\omega(q)}m_r(q)}\right)^{-1/r} \notag \\
    %&=
    %\frac{2^{1+1/r}(r!)^{1/r}}{r^{1-1/r}(r-1)^{1+1/r}\log q} \bigl(\left(4r\right)^{\omega(q)}m_r(q)\bigr)^{1+1/r}.
\end{equation}

\begin{theorem}\label{main thm 1}
Let $r\ge 2$ be an integer and $\chi$ be a primitive Dirichlet character modulo~$q$. Let $C(r)$ be defined as in Table \ref{Table 1}. Let $a(r) = 2\log{2}\left(3.0758r+1.38402\log(4r)-1.5379\right)$. Then, for $q \ge \max\{10^{1143},e^{e^{a(r)}}\}$, if $r=2$ or $q$ is cubefree, we have
    $$|S_{\chi}(M,N)| \le C(r)N^{1-\frac{1}{r}} q^{\frac{r+1}{4r^2}}(\log{q})^{\frac{1}{2r}}\left((4r)^{\omega(q)}m_r(q)\right)^{\frac{1}{2r}-\frac{1}{2r^2}}\left(\frac{q}{\phi(q)}\right)^{\frac{1}{r}}.$$

    Furthermore, as $q\rightarrow\infty$, we have a constant $D(r)$ from Table \ref{Table 1} such that
    $$|S_{\chi}(M,N)| \le (D(r)+o(1))N^{1-\frac{1}{r}} q^{\frac{r+1}{4r^2}}(\log{q})^{\frac{1}{2r}}\left((4r)^{\omega(q)}m_r(q)\right)^{\frac{1}{2r}-\frac{1}{2r^2}}\left(\frac{q}{\phi(q)}\right)^{\frac{1}{r}}.$$
\end{theorem}

\begin{proof}[Proof of Theorem \ref{main thm 1}]
Recall that we assume $q\ge \max\{10^{1143}, t_1(r)\}$ for $t_1(r)=e^{e^{a(r)}}$ with $a(r)$ given in the statement of Theorem \ref{main thm 1} (and Lemma \ref{B will be at least 10 if it kills us lemma}), and from
\eqref{eq T1}
$$T_1(q) = \left((4r)^{\omega(q)}m_r(q)\right)^{\frac{1}{2r}-\frac{1}{2r^2}}\left(\frac{q}{\phi(q)}\right)^{\frac{1}{r}}.$$
By Lemma \ref{lemma to set up the main theorems}
and comparing with \eqref{eq desired bound}, it suffices to show that
\begin{equation*}
u(\alpha w)^{\frac{1}{2r}} + sT_1(q) \le \frac{B-1}{B} C T_1(q).
\end{equation*}

Since $10^{1143}\ge (10r)^{18}$ for $r\le 10^{50}$ while $t_1(r) \ge (10r)^{18}$ for $r\ge 10^{50}$, we apply Lemma \ref{A lower bound lemma} and conclude that $A\phi^{*} \ge 2^{\omega(q)}$, hence
$$\frac{N}{\#\mathcal{A}}\le \frac{N}{A\phi^*-2^{\omega(q)-1}}=\frac{N}{A\phi^*}+\frac{2^{\omega(q)-1}N}{A\phi^*(A\phi^*-2^{\omega(q)-1})} \le \frac{2N}{A\phi^{*}}.$$
Therefore, \eqref{eq u} satisfies the bound
\begin{equation} \label{roughly u}
u=\frac{\left(N/\#\cA\right)^{1/{r}}}{q^\frac{1}{2r^2}} \le \frac{2^{1/r}}{q^{\frac{1}{2r^2}}}\left(\frac{N}{A\phi^*}\right)^{1/r}.
\end{equation}

From Lemma \ref{lem: alpha bound}, we have
\begin{align}\label{eq alpha upper bound}
\alpha \le \frac{32}{31}\kappa \log{q}.    
\end{align}
Using that $f_1 = r$ and \begin{equation}\label{beta B def}\beta = \left(\frac{r-1}{(2r)r!(4r)^{\omega(q)}m_r(q)}\right)^{\frac{1}{r}},
\end{equation}
we rewrite \eqref{w} as
\begin{equation}\label{bounding w}
    w = \frac{2^{1+1/r}(r\cdot r!)^{1/r}}{(r-1)^{1+1/r}\log q} \bigl(\left(4r\right)^{\omega(q)}m_r(q)\bigr)^{1+1/r}.
\end{equation}
Furthermore, using \eqref{beta B def}, $B = r^2q^{\frac{1}{2r}}\beta$, and $AB = \kappa N$, we have
\begin{equation}\label{bounding N over A}
\frac{N}{A q^{\frac{1}{2r}}}\left((4r)^{\omega(q)}m_r(q)\right)^{\frac{1}{r}} = \frac{r^2}{\kappa}\left(\frac{r-1}{r!(2r)}\right)^{\frac{1}{r}}.
\end{equation}
Therefore, combining \eqref{eq s}, \eqref{roughly u}, \eqref{eq alpha upper bound}, \eqref{bounding w}, and \eqref{bounding N over A},  we obtain
\begin{align*}
    &u(\alpha w)^{\frac{1}{2r}}+sT_1(q)\\ 
    &\le 2^{\frac{1}{r}}\left(\frac{N}{Aq^{\frac{1}{2r}}\phi^*}\right)^{\frac{1}{r}}\left(\frac{32}{31}\kappa \log{q}\right)^{\frac{1}{2r}}\left(\frac{2^{1+\frac{1}{r}}(r\cdot r!)^{\frac{1}{r}}}{(r-1)^{1+\frac{1}{r}}\log{q}}\right)^{\frac{1}{2r}}\left((4r)^{\omega(q)}m_r(q)\right)^{\frac{1}{2r}+\frac{1}{2r^2}}\\ &+\frac{2C\kappa^{1-\frac{1}{r}}}{2-\frac{1}{r}}\left(\frac{B+1}{B}\right)^{2-\frac{1}{r}}T_1(q)\\
    &= \left(2^{\frac{4}{r}-\frac{1}{2r^2}}31^{-\frac{1}{2r}}(r!)^{-\frac{1}{2r^2}}(r-1)^{\frac{1}{2r^2}-\frac{1}{2r}}r^{\frac{2}{r}-\frac{1}{2r^2}}\kappa^{-\frac{1}{2r}} +\frac{2C\kappa^{1-\frac{1}{r}}}{2-\frac{1}{r}}\left(\frac{B+1}{B}\right)^{2-\frac{1}{r}}\right)T_1(q).
\end{align*}
Therefore,
\begin{equation} \label{C lower bound}
C \ge \frac{2^{\frac{4}{r}-\frac{1}{2r^2}}31^{-\frac{1}{2r}}(r!)^{-\frac{1}{2r^2}}(r-1)^{\frac{1}{2r^2}-\frac{1}{2r}}r^{\frac{2}{r}-\frac{1}{2r^2}}\kappa^{-\frac{1}{2r}}}{\frac{B-1}{B}-\frac{2\kappa^{1-\frac{1}{r}}}{2-\frac{1}{r}}\left(\frac{B+1}{B}\right)^{2-\frac{1}{r}}}.
\end{equation}

Now to obtain the constants in Table \ref{Table 1}, first, by Lemmas \ref{B will be at least 10 if it kills us lemma} and \ref{B will be at least 10 if it kills us lemma 2},
\begin{equation}\label{B lower bound123}
B\ge 2^{2-\frac{2}{r}}r^2\left(\frac{1-\frac{1}{r}}{r!}\right)^{\frac{1}{r}}.
\end{equation}
Then, choosing $\kappa$ as in \eqref{choice of kappa}, we replace $B$ with $2^{2-2/r}r^2\left(\frac{1-\frac{1}{r}}{r!}\right)^{\frac{1}{r}}$ in $\kappa$ since the right side of \eqref{C lower bound} is decreasing as $B$ increases for $B\ge 2^{2-2/r}r^2\left(\frac{1-\frac{1}{r}}{r!}\right)^{\frac{1}{r}}$. Computing the resulting expression obtains the $C(r)$ column in Table \ref{Table 1}. 
For the $D(r)$ column in Table \ref{Table 1}, we notice that $q\rightarrow \infty$ implies $B\rightarrow \infty$, therefore $$\kappa\rightarrow \left(\frac{1}{2r}\right)^{\frac{r}{r-1}},$$
and we obtain
$$D(r) = \frac{2^{\frac{4}{r}-\frac{1}{2r^2}}31^{-\frac{1}{2r}}(r!)^{-\frac{1}{2r^2}}(r-1)^{\frac{1}{2r^2}-\frac{1}{2r}}r^{\frac{2}{r}-\frac{1}{2r^2}}(2r)^{\frac{1}{2r-2}}}{1-\frac{1}{2r-1}}.$$
\end{proof}

\end{document}
